\section{Limitations}

The present solver is intentionally reduced-order. It does not solve full transient Navier--Stokes equations, full Taylor--Melcher leaky-dielectric surface charge conservation, droplet breakup, plume dynamics, or complete thruster performance. The effective space-charge closures are phenomenological models for studying Poisson shielding trends, not predictive emitted-current laws.

The axisymmetric assumption is appropriate for onset-level symmetric geometries but may fail after non-axisymmetric instabilities or downstream jet breakup appear. The finite-difference grid representation introduces geometric error, especially near sharp conical interfaces. The refinement study quantifies this: the recovered half-angle approaches the ideal Taylor value as the grid is refined ($\SI{49.14}{\degree}$, $\SI{49.20}{\degree}$ and $\SI{49.23}{\degree}$ at the three tested levels), so the remaining deviation is a measured discretization error rather than an unquantified one.

Experimental validation will also have limitations. Cone-angle extraction depends on image resolution, optical alignment, lighting, threshold selection, calibration accuracy, and frame selection. Onset voltage may depend on fluid impurities, humidity, flow stability, electrode condition, and subjective stability criteria unless carefully standardized. Therefore, final conclusions will be stated in terms of quantified agreement and uncertainty rather than exact prediction.
